\documentclass[11pt,a4paper]{article}
\usepackage[T1]{fontenc}
\usepackage[utf8]{inputenc}
\usepackage[margin=2.3cm]{geometry}
\usepackage{lmodern}
\usepackage{microtype}
\usepackage{booktabs,array,tabularx}
\usepackage{eurosym}
\usepackage[hidelinks]{hyperref}
\setlength{\parindent}{0pt}
\setlength{\parskip}{0.5em}
\newcommand{\keur}{\EUR\,k}
\title{ARGOS: initial investment estimate}
\author{Working budget for discussion}
\date{13 September 2026}
\begin{document}
\maketitle

\section{Scope and proposed funding}

The proposed programme builds one PET scanner, one Compton camera and the
integrated software and service needed for a full system demonstrator.
Development with Donostia Hospital proceeds from phantom tests to clinical
investigations, clinical validation and conformity assessment.

The central estimate is \textbf{\EUR\,6.7 million over four years}, with a
\textbf{funding target of \EUR\,7 million}. A first investment of
\EUR\,2--3 million could finance the initial stage, but would not cover the
complete programme and company operations.

The estimate assumes a Donostia-based team, access to hospital and proton-beam
facilities by agreement, and no construction of a proton-therapy facility.
Costs are planning allowances in 2026 euros, excluding recoverable VAT.
No sales revenue, grants or partner contributions are deducted. Price increases
and schedule overruns draw on the programme reserve.

\section{Four-year budget}

\begin{table}[htbp]
\centering\small
\renewcommand{\arraystretch}{1.15}
\begin{tabularx}{\textwidth}{@{}p{4.5cm}rX@{}}
\toprule
Item & \keur & Scope \\
\midrule
PET, Compton camera and computing & 1,000 & Existing raw-hardware estimate. \\
Prototype fabrication and test equipment & 200 & Purchased fabrication, phantoms, fixtures, installation and external technical work. \\
First-device contingency & 300 & Failed components, redesigns and replacement parts. \\
Personnel & 2,660 & Development, management, quality and clinical monitoring. \\
Hospital and clinical programme & 500 & Phantom sessions, beam access, hospital staff, trial coordination, data management and study insurance. \\
External regulatory and conformity work & 250 & Specialist advice, independent testing, audits and certification work. \\
Office and laboratory space & 240 & \EUR\,60k/year, including utilities. \\
Software tools and IT operations & 120 & \EUR\,30k/year: tools, security, backup and licences. \\
Company administration & 160 & \EUR\,40k/year: accounting, legal work, IP and corporate insurance. \\
Travel and transport & 120 & \EUR\,30k/year: hospital visits, suppliers, conferences and equipment transport. \\
\midrule
\textbf{Planned expenditure} & \textbf{5,550} & Includes the first-device contingency. \\
Programme reserve & 810 & 20\% of the \EUR\,4.05M outside the hardware/prototype package. \\
Working-capital buffer & 300 & Payment timing, recoverable VAT and cash at programme completion. \\
\midrule
\textbf{Funding requirement} & \textbf{6,660} & \textbf{Round to \EUR\,6.7M; target \EUR\,7M.} \\
\bottomrule
\end{tabularx}
\end{table}

The raw-hardware estimate comprises \EUR\,740k for PET materials,
\EUR\,100k for the Compton camera and \EUR\,160k for computing and other
hardware. These inputs follow the existing installation model.

The \EUR\,200k prototype allowance covers external purchases and contracted
work. Employee assembly, integration and software development are charged to
personnel. The full delivered-system cost is not added again. The first-device
contingency is 25\% of the \EUR\,1.2M hardware and external-work base; the
programme reserve applies only to the remaining expenditure.

Working capital is a cash requirement, not an additional development expense.
The difference between \EUR\,6.66M and the \EUR\,7M funding target provides
\EUR\,340k of further headroom.

\section{Personnel}

Annual figures are total employer costs, including employer contributions,
not gross salaries. They are proposed hiring allowances rather than salary
survey results. Founders taking operating roles are included in the paid team.

\begin{table}[htbp]
\centering\small
\renewcommand{\arraystretch}{1.2}
\begin{tabularx}{\textwidth}{@{}Xlrr@{}}
\toprule
Role & Staffing period & Annual \keur & Total \keur \\
\midrule
CEO / project and business lead & $1\times4$ years & 90 & 360 \\
System architect / technical lead & $1\times4$ years & 100 & 400 \\
Detector and electronics engineer & $1\times4$ years & 85 & 340 \\
Software engineers & $2\times4$ years & 85 each & 680 \\
Medical physicist / validation lead & $1\times4$ years & 90 & 360 \\
Assembly and laboratory technician & $1\times3$ years & 60 & 180 \\
Quality and regulatory lead & $0.5\times4$ years & 90 per FTE & 180 \\
Clinical monitoring specialist & $1\times2$ years & 80 & 160 \\
\midrule
\textbf{Total} & & & \textbf{2,660} \\
\bottomrule
\end{tabularx}
\end{table}

The software engineers develop reconstruction, data integration, interfaces
and service tools. Their costs are company development expenditure; they are
not included in the non-personnel IT allowance.

The medical physicist leads validation. A separate monitoring specialist
operates the daily clinical service during the two clinical years. The
\EUR\,80k/year provision follows the existing monitoring-cost assumption,
including relief coverage for one shift. Hospital staff and trial costs are
in the separate \EUR\,500k clinical allowance, which should be replaced by a
costed hospital agreement once the protocol and patient numbers are defined.

\section{Development sequence}

\begin{description}
\item[Year 1.] Detailed design, procurement, software architecture, quality
system and clinical protocol.
\item[Year 2.] Integrated prototype, phantom measurements, reliability testing
and service procedures.
\item[Years 3--4.] Authorised clinical investigations, system improvements,
clinical validation and conformity assessment.
\end{description}

Donostia Hospital is the clinical development and validation partner. Device
market certification follows the applicable conformity-assessment route;
the hospital does not itself grant it. Spanish clinical investigations require
the relevant ethics review and AEMPS authorisation~\cite{aemps}; notified bodies
assess conformity where required by the device's regulatory route~\cite{ec}.

The funded endpoint is a working, clinically evaluated demonstrator and its
service, with European conformity work included. The budget does not cover a
large multicentre study proving improved patient outcomes, worldwide market
approvals or production facilities for commercial-scale manufacturing.

\section{Effect of programme duration}

\begin{center}
\begin{tabular}{@{}lr@{}}
\toprule
Development period & Funding range \\
\midrule
3 years & \EUR\,5.5--6M \\
\textbf{4 years} & \textbf{\EUR\,6.5--7M} \\
5 years & \EUR\,7.5--8M \\
\bottomrule
\end{tabular}
\end{center}

These are indicative duration variants, not separately costed programmes.
The three-year case requires substantial design maturity and prompt access
to clinical testing; compressing the schedule does not automatically reduce
engineering effort. Retaining the team and facilities for an additional year
costs approximately \EUR\,0.8--0.9M before contingency.

The proposed starting point is \textbf{\EUR\,7M over four years}, potentially
raised in stages against prototype and clinical milestones. This note records
the budget proposal separately from the six existing financial and market
tables; it does not alter the presentation or workbook.

\begin{thebibliography}{9}
\small
\bibitem{aemps} AEMPS, \emph{Investigaciones cl\'inicas con productos sanitarios}.
\url{https://www.aemps.gob.es/productos-sanitarios/investigacionclinica-productossanitarios/}.
\bibitem{ec} European Commission, \emph{Notified bodies for medical devices}.
\url{https://health.ec.europa.eu/medical-devices-topics-interest/notified-bodies-medical-devices_en}.
\end{thebibliography}
\end{document}
